\chapter{Машина целиком}
% полная версия: chapters/10_synthesis.tex

\pencilfig{10}{\pencilart{10}}{Машина целиком: шина данных, тормоза в ней и четыре радиоканала управления сверху.}

Мы брали по одному типу нейронов и спрашивали, что он добавляет к машине. Пора собрать детали вместе.

Получилась машина из трёх слоёв. Внизу --- огромная телефонная сеть \nt{ГЛУТ}-нейронов: по ней идут данные. Рядом --- тормозные \nt{ГАМК}-нейроны: они управляют потоком, решают, чему пройти, когда и насколько громко. А над всем этим --- несколько радиостанций, у каждой своя ручка. \nt{ДОФА} говорит, какие изменения связей оставить. \nt{СЕРО} держит общий уровень и, по гипотезе, горизонт терпения. \nt{НОРА} выбирает между поиском и решением. \nt{АЦЕТ} --- между записью и чтением. Каждая радиостанция на самом деле --- небольшой пучок каналов, а не один провод, но каналов всё равно горстка на восемьдесят шесть миллиардов нейронов.

\section*{Одна история}

Проследим одно событие через всю машину. Животное давно знает: нажмёшь рычаг A --- получишь сок. Однажды мир меняется: A больше не работает, а сок приносит рычаг B, которым животное почти не пользовалось.

Сок после A не пришёл --- дофамин отвечает провалом, и связи, выбравшие A, слабеют. Провалы повторяются, ошибки становятся больше обычного --- норадреналин, по гипотезе, сообщает: мир изменился. Ручка контрастности опускается, выбор становится мягче, и шум чаще подталкивает попробовать B. Ацетилхолин переводит сеть в режим записи. Проба B приносит неожиданный сок --- дофамин вспыхивает, связи, выбравшие B, крепнут. Через несколько попыток ожидание догоняет новый мир, сюрпризы кончаются, ручки возвращаются на место.

Заметьте: ни одна ручка не знает, что делают другие. Каждая откликается на свои признаки, и порядок действий складывается сам, как в асинхронной схеме, где каждый блок срабатывает, когда готовы его входы. Что каждая ручка ведёт себя так по отдельности, в основном измерено. Что они слаженно работают вместе --- пока гипотеза.

\section*{Чем она не похожа на ваш ноутбук}

У ноутбука есть общий такт --- у мозга нет. Элементы ноутбука надёжны --- синапс теряет больше половины сигналов. Память ноутбука отдельно от процессора --- в мозге она в тех же связях. Ноутбук учится (если учится) обратным распространением ошибки --- мозг местными правилами и одним широковещательным «лучше, чем ожидалось». И главное: ноутбук тратит десятки ватт на один процессор, а мозг --- двадцать ватт на всё.

\section*{Чего мы не знаем}

Мы не знаем, насколько кора пользуется точным временем щелчков. Не знаем, как она настраивает миллиарды связей в многоэтажных цепях, если учитель у неё один на всех. Не до конца понимаем, что передают радиостанции. Не знаем, как далёкие части мозга согласуют работу без общего такта. И пока никто не умеет построить машину, которая делала бы то же самое на двадцати ваттах.

Дальше книга выходит за пределы этого ядра: к входам и выходам машины, к её отладке, ко сну и к машинам из живых нейронов на чипе.

\fullversion{10}
